% Options for packages loaded elsewhere
\PassOptionsToPackage{unicode}{hyperref}
\PassOptionsToPackage{hyphens}{url}
\documentclass[
]{article}
\usepackage{xcolor}
\usepackage{amsmath,amssymb}
\setcounter{secnumdepth}{-\maxdimen} % remove section numbering
\usepackage{iftex}
\ifPDFTeX
  \usepackage[T1]{fontenc}
  \usepackage[utf8]{inputenc}
  \usepackage{textcomp} % provide euro and other symbols
\else % if luatex or xetex
  \usepackage{unicode-math} % this also loads fontspec
  \defaultfontfeatures{Scale=MatchLowercase}
  \defaultfontfeatures[\rmfamily]{Ligatures=TeX,Scale=1}
\fi
\usepackage{lmodern}
\ifPDFTeX\else
  % xetex/luatex font selection
\fi
% Use upquote if available, for straight quotes in verbatim environments
\IfFileExists{upquote.sty}{\usepackage{upquote}}{}
\IfFileExists{microtype.sty}{% use microtype if available
  \usepackage[]{microtype}
  \UseMicrotypeSet[protrusion]{basicmath} % disable protrusion for tt fonts
}{}
\makeatletter
\@ifundefined{KOMAClassName}{% if non-KOMA class
  \IfFileExists{parskip.sty}{%
    \usepackage{parskip}
  }{% else
    \setlength{\parindent}{0pt}
    \setlength{\parskip}{6pt plus 2pt minus 1pt}}
}{% if KOMA class
  \KOMAoptions{parskip=half}}
\makeatother
\usepackage{longtable,booktabs,array}
\usepackage{calc} % for calculating minipage widths
% Correct order of tables after \paragraph or \subparagraph
\usepackage{etoolbox}
\makeatletter
\patchcmd\longtable{\par}{\if@noskipsec\mbox{}\fi\par}{}{}
\makeatother
% Allow footnotes in longtable head/foot
\IfFileExists{footnotehyper.sty}{\usepackage{footnotehyper}}{\usepackage{footnote}}
\makesavenoteenv{longtable}
\setlength{\emergencystretch}{3em} % prevent overfull lines
\providecommand{\tightlist}{%
  \setlength{\itemsep}{0pt}\setlength{\parskip}{0pt}}
\usepackage{bookmark}
\IfFileExists{xurl.sty}{\usepackage{xurl}}{} % add URL line breaks if available
\urlstyle{same}
\hypersetup{
  hidelinks,
  pdfcreator={LaTeX via pandoc}}

\author{}
\date{}

\begin{document}

\section{\texorpdfstring{Universal Obstruction Theory: The \(\ell^1\)
Topological
Framework}{Universal Obstruction Theory: The \textbackslash ell\^{}1 Topological Framework}}\label{universal-obstruction-theory-the-ell1-topological-framework}

\textbf{Author:} Jeremy H. Carroll \textbf{Date:} March 2026
\textbf{Version:} v003 (Pre-print)

\begin{center}\rule{0.5\linewidth}{0.5pt}\end{center}

\subsection{Classification of Results}\label{classification-of-results}

\begin{longtable}[]{@{}
  >{\raggedright\arraybackslash}p{(\linewidth - 4\tabcolsep) * \real{0.2121}}
  >{\raggedright\arraybackslash}p{(\linewidth - 4\tabcolsep) * \real{0.2424}}
  >{\raggedright\arraybackslash}p{(\linewidth - 4\tabcolsep) * \real{0.5455}}@{}}
\toprule\noalign{}
\begin{minipage}[b]{\linewidth}\raggedright
Layer
\end{minipage} & \begin{minipage}[b]{\linewidth}\raggedright
Status
\end{minipage} & \begin{minipage}[b]{\linewidth}\raggedright
What it contains
\end{minipage} \\
\midrule\noalign{}
\endhead
\bottomrule\noalign{}
\endlastfoot
\textbf{A: Established} & Proved in Papers 000--003 & \(\ell^1\)
classification (combinatorial, Banach, dynamical), Hodge quotient norm,
density persistence \\
\textbf{B: Theorems} & Proved in this paper & Functorial aggregation
monotonicity, operator-valued extension to Schatten-1, sheaf Laplacian
gradient structure \\
\textbf{C: Conjectural} & Research program & Born rule emergence from
\(\ell^1\) dynamics, causal set extensions \\
\end{longtable}

\begin{center}\rule{0.5\linewidth}{0.5pt}\end{center}

\subsection{Abstract}\label{abstract}

Papers 000--003 established, through independent proofs at four
structural levels, that any additive, faithful, structure-preserving
diagnostic on coboundary defects is uniquely forced to be the \(\ell^1\)
coboundary norm. Paper 002 proved this at the combinatorial level
(coordinate additivity on finite graph cochains). Paper 001 extended it
to Banach presheaves (functoriality under bounded linear maps). Paper
000 showed that projection dynamics generate intrinsic defects measured
by this norm, persisting at non-vanishing density. Paper 003
demonstrated that the forced \(\ell^1\) geometry induces a canonical
gauge-invariant quotient norm on first cohomology.

This paper synthesizes these results into a unified obstruction
framework and extends them in three directions. First, we prove that the
\(\ell^1\) coboundary norm is monotone under coarse-graining functors:
topological defects cannot vanish upon aggregation unless their local
boundaries physically cancel (Theorem 2.1). Second, we extend the
framework from scalar and vector-valued presheaves to operator-valued
presheaves, showing that when stalks are spaces of bounded operators on
Hilbert spaces, the \(\ell^1\) coboundary norm lifts canonically to the
Schatten-1 (trace) norm (Theorem 3.1). Third, we identify the sheaf
Laplacian as the natural gradient flow operator for \(\ell^1\) defect
minimization and characterize its fixed points as globally consistent
sections (Theorem 4.1).

We further outline, at the level of a research program, how these three
extensions connect the foundational quad to subsequent applications in
autopoietic cohomology, gauge group emergence, the fermion sign problem,
and the behavior of coupling constants. All conjectural claims are
explicitly marked.

\begin{center}\rule{0.5\linewidth}{0.5pt}\end{center}

\subsection{1. Introduction}\label{introduction}

\subsubsection{1.1 The Foundational Quad}\label{the-foundational-quad}

The preceding four papers establish a single structural conclusion at
four levels of generality:

\begin{enumerate}
\def\labelenumi{\arabic{enumi}.}
\item
  \textbf{Combinatorial (Paper 002)} {[}6{]}: On finite graph cochains,
  edge locality, coordinate additivity over disjoint supports,
  coordinate symmetry, and graph invariance force the \(\ell^1\) norm as
  the unique diagnostic geometry, up to a global scalar.
\item
  \textbf{Functional-analytic (Paper 001)} {[}5{]}: On Banach presheaf
  coboundaries, edge-wise additivity, faithfulness, functoriality under
  bounded linear endomorphisms, and isomorphism invariance force the
  \(\ell^1\) coboundary norm as the unique edge-additive seminorm, up to
  a positive scalar.
\item
  \textbf{Dynamical (Paper 000)} {[}4{]}: Projection of linear many-body
  evolution onto a nonlinear manifold produces intrinsic coboundary
  defects that (a) obstruct linear conjugacy of the projected flow, (b)
  are measured by the uniquely forced \(\ell^1\) diagnostic, and (c)
  persist at non-vanishing density over a finite time window independent
  of system size.
\item
  \textbf{Cohomological (Paper 003)} {[}7{]}: The forced \(\ell^1\)
  geometry on \(C^1(G)\) induces a canonical quotient norm \(\Phi\) on
  \(H^1(G)\), providing the unique gauge-invariant measure of
  irreducible obstruction. The \(\ell^2\) Hodge projection generically
  inflates \(\ell^1\) cost by a factor approaching 3 on cycle graphs.
\end{enumerate}

The present paper is not a new proof of \(\ell^1\) rigidity. It is a
synthesis: we organize the quad into a single categorical framework,
extend it to new structural settings, and identify the connections to
subsequent work.

\subsubsection{1.2 What This Paper
Contributes}\label{what-this-paper-contributes}

Three theorems extend the quad:

\begin{itemize}
\item
  \textbf{Theorem 2.1 (Functorial Aggregation Monotonicity):} The
  \(\ell^1\) coboundary norm is monotonically non-increasing under
  coarse-graining functors. Defects cannot increase under aggregation,
  and can only vanish if locally canceled.
\item
  \textbf{Theorem 3.1 (Operator-Valued Extension):} When presheaf stalks
  are operator algebras \(\mathcal{B}(\mathcal{H})\), the \(\ell^1\)
  coboundary norm lifts to the Schatten-1 (trace) norm, preserving all
  four classification axioms.
\item
  \textbf{Theorem 4.1 (Sheaf Laplacian Gradient Structure):} The sheaf
  Laplacian \(\Delta_\mathcal{F}\) governs the gradient flow of the
  squared \(\ell^2\) defect functional on sections, and its kernel
  characterizes globally consistent sections.
\end{itemize}

\subsubsection{1.3 Relation to the Full
Program}\label{relation-to-the-full-program}

The program proceeds in stages:

\begin{itemize}
\item
  Papers 000--003: Establish the \(\ell^1\) obstruction framework across
  combinatorial, functional, dynamical, and cohomological settings.
\item
  Paper 004 (this work): Synthesizes and extends the framework
  (aggregation, operator-valued extension, dynamical structure).
\item
  Subsequent papers: Explore applications of the framework to
  increasingly structured physical and computational settings, including
  gauge structure, fermionic systems, and coupling behavior.
\end{itemize}

\subsubsection{1.4 Computational
Companion}\label{computational-companion}

This paper is accompanied by an illustrative Python execution script
(\texttt{universal\_obstruction\_derivation.py}) that models the
sequential behavior of Papers 000--004. It demonstrates topological
defect resolution via zero-crossing subdivision, calculates discrete
Banach seminorms, and provides a direct contrast between \(\ell^1\) cost
minimization and classical \(\ell^2\) Hodge decomposition. This script
serves as a strictly structural illustration of the mathematical
framework rather than a formal physical simulation.

\subsubsection{1.5 Conventions}\label{conventions}

Norms are denoted \(\lVert \cdot \rVert\) with appropriate subscripts.
The Schatten-\(p\) norm on \(\mathcal{B}(\mathcal{H})\) is
\(\|A\|_p = (\mathrm{Tr}|A|^p)^{1/p}\), where \(|A| = (A^*A)^{1/2}\).
The Schatten-1 norm is the trace norm \(\|A\|_1 = \mathrm{Tr}|A|\).
Scalar absolute values are denoted \(|\cdot|\).

\begin{center}\rule{0.5\linewidth}{0.5pt}\end{center}

\subsection{2. Functorial Defect
Aggregation}\label{functorial-defect-aggregation}

\subsubsection{2.1 Setup}\label{setup}

Let \((P, \leq)\) be a finite poset with covering relations
\(E_{\mathrm{cov}}\), and let \(F\) be a Banach presheaf over \(P\). A
\textbf{coarse-graining functor} is a surjective order-preserving map
\(\pi: P \to Q\) to a smaller poset \(Q\), together with a natural
transformation that aggregates the stalks: for each \(q \in Q\), the
aggregated stalk is \(G(q) = \bigoplus_{x \in \pi^{-1}(q)} F(x)\) with
restriction maps induced from \(F\).

The \textbf{pushforward} of a section \(s\) of \(F\) is the section
\(\pi_* s\) of \(G\) defined by
\((\pi_* s)(q) = (s(x))_{x \in \pi^{-1}(q)}\).

\subsubsection{2.2 Monotonicity}\label{monotonicity}

\begin{quote}
\textbf{Theorem 2.1 (Functorial Aggregation Monotonicity).}
\textbf{(Layer B)} Let \(\pi: (P, F) \to (Q, G)\) be a coarse-graining
as above. Then \[I_Q(\pi_* s) \leq I_P(s)\] for every section \(s\) of
\(F\). Equality holds if and only if every edge of \(P\) that is
collapsed by \(\pi\) (i.e., both endpoints map to the same element of
\(Q\)) has zero coboundary defect \(\delta_e(s) = 0\).
\end{quote}

\emph{Proof.} The \(\ell^1\) coboundary norm on \(P\) sums over all
covering edges of \(P\):
\[I_P(s) = \sum_{e \in E_{\mathrm{cov}}(P)} \|\delta_e(s)\|\]

Under \(\pi\), the edges of \(P\) partition into two classes: (a) edges
\(e = (a, b)\) with \(\pi(a) \neq \pi(b)\), which map to covering edges
of \(Q\), and (b) edges \(e = (a, b)\) with \(\pi(a) = \pi(b)\), which
are collapsed.

For edges of type (a), the coboundary defect on \(Q\) satisfies
\(\|\delta_{\pi(e)}(\pi_* s)\| \leq \|\delta_e(s)\|\) by contractivity
of the induced restriction maps on the aggregated stalks.

Edges of type (b) contribute to \(I_P(s)\) but not to \(I_Q(\pi_* s)\),
since they are no longer covering relations in \(Q\).

Therefore:
\[I_Q(\pi_* s) \leq \sum_{e: \pi(a) \neq \pi(b)} \|\delta_e(s)\| \leq I_P(s)\]

Equality requires both that the type-(a) contractions are isometric and
that all type-(b) edges have zero defect. \(\square\)

\textbf{Interpretation.} Topological defects cannot be hidden by
coarse-graining. An \(\ell^1\) defect that is present at the fine scale
either persists at the coarse scale or was already zero on the collapsed
edges. This is the structural reason why \(\ell^1\) obstruction theory
is stable under renormalization: the defect measure is monotone, not
merely invariant.

\begin{center}\rule{0.5\linewidth}{0.5pt}\end{center}

\subsection{3. Operator-Valued Presheaves and the Schatten-1
Norm}\label{operator-valued-presheaves-and-the-schatten-1-norm}

\subsubsection{3.1 From Scalars to
Operators}\label{from-scalars-to-operators}

Papers 001 and 002 classify the \(\ell^1\) norm on presheaves with
scalar or vector-valued stalks. In quantum many-body systems, the
natural stalks are spaces of density operators on local Hilbert spaces.
This section extends the classification to this setting.

Let \((P, \leq)\) be a finite poset. An \textbf{operator-valued
presheaf} \(F\) over \(P\) assigns: - To each \(x \in P\), a
finite-dimensional Hilbert space \(\mathcal{H}_x\) and the Banach space
\(F(x) = \mathcal{B}_1(\mathcal{H}_x)\) of trace-class operators,
equipped with the Schatten-1 (trace) norm \(\|A\|_1 = \mathrm{Tr}|A|\).
- To each covering edge \(e = (a, b)\), a completely positive
trace-non-increasing restriction map \(r_{b,a}: F(b) \to F(a)\) (the
partial trace or a conditional expectation).

The coboundary defect at edge \(e = (a, b)\) for a section
\(\rho = (\rho_x)_{x \in P}\) (where \(\rho_x \in F(x)\) is a local
density operator) is:
\[\delta_e(\rho) = r_{b,a}(\rho_b) - \rho_a \in F(a)\]

This measures the failure of the local state \(\rho_a\) to be the
marginal of \(\rho_b\) --- precisely the quantum marginal
incompatibility.

\subsubsection{3.2 Classification
Extension}\label{classification-extension}

\begin{quote}
\textbf{Theorem 3.1 (Operator-Valued \(\ell^1\) Classification).}
\textbf{(Layer B)} Let \(\nu\) be a seminorm on the operator-valued
1-cochain space
\(C^1(P, F) = \bigoplus_{e \in E_{\mathrm{cov}}} \mathcal{B}_1(\mathcal{H}_{a_e})\)
satisfying:

(P1) Edge-wise additivity:
\(\nu((\delta_e)_e) = \sum_e \nu_e(\delta_e)\)

(P2) Faithfulness: \(\nu_e(\delta) = 0 \iff \delta = 0\)

(P3) Functoriality:
\(\nu_e(\Phi(\delta)) \leq \|\Phi\|_{\mathrm{cb}} \cdot \nu_e(\delta)\)
for every completely bounded map
\(\Phi: \mathcal{B}_1(\mathcal{H}) \to \mathcal{B}_1(\mathcal{H})\)

(P4) Isomorphism invariance

Then \(\nu_e(\delta) = w_e \cdot \|\delta\|_1\) for some \(w_e > 0\),
where \(\|\cdot\|_1\) is the Schatten-1 (trace) norm. If (P4) holds,
\(w_e = \alpha\) for all \(e\).
\end{quote}

\emph{Proof sketch.} The argument follows the structure of Paper 001,
Theorem 2.2. The key step is the direction-independence lemma: for any
two operators \(A, B \in \mathcal{B}_1(\mathcal{H})\) with
\(\|A\|_1 = \|B\|_1 \neq 0\), we must show \(\nu_e(A) = \nu_e(B)\).

By the operator-valued Hahn--Banach theorem (a consequence of the
duality \(\mathcal{B}_1(\mathcal{H})^* = \mathcal{B}(\mathcal{H})\)),
for any \(A\) with \(\|A\|_1 = 1\), there exists
\(X \in \mathcal{B}(\mathcal{H})\) with \(\|X\|_{\mathrm{op}} = 1\) and
\(\mathrm{Tr}(X^* A) = 1\). Define the completely bounded map
\(\Phi: \mathcal{B}_1(\mathcal{H}) \to \mathcal{B}_1(\mathcal{H})\) by
\(\Phi(C) = \mathrm{Tr}(X^* C) \cdot B / \|B\|_1\). Then
\(\Phi(A) = B\), \(\|\Phi\|_{\mathrm{cb}} = 1\), and (P3) gives
\(\nu_e(B) \leq \nu_e(A)\). By symmetry, equality follows.

The remainder proceeds identically to the scalar case: direction
independence plus homogeneity forces \(\nu_e = w_e \|\cdot\|_1\), and
(P4) forces uniform weights. \(\square\)

\subsubsection{3.3 Structural
Consequences}\label{structural-consequences}

\textbf{Quantum marginal incompatibility.} When \(\rho\) is a global
quantum state and
\(\Pi(\rho) = \bigotimes_x \mathrm{Tr}_{\bar{x}}(\rho)\) is the
product-state retraction, the coboundary defect \(\delta_e(\Pi(\rho))\)
measures the entanglement across edge \(e\) in trace distance. Theorem
3.1 shows this is the unique additive, faithful, functorial measure of
marginal incompatibility.

\textbf{Entanglement monogamy.} The edge-wise additivity (P1) combined
with the trace norm structure implies that total defect across all edges
sharing a vertex is bounded by the local dimension. This recovers the
structure of entanglement monogamy inequalities as a consequence of the
\(\ell^1\) framework, rather than as an independent constraint.

\begin{center}\rule{0.5\linewidth}{0.5pt}\end{center}

\subsection{4. Dynamic Selection via the Sheaf
Laplacian}\label{dynamic-selection-via-the-sheaf-laplacian}

\subsubsection{4.1 The Sheaf Laplacian}\label{the-sheaf-laplacian}

Given a Banach presheaf \(F\) over a finite poset \(P\), the
\textbf{coboundary operator} \(d: C^0(P, F) \to C^1(P, F)\) maps
sections to their edge defects:
\((ds)_e = \delta_e(s) = r_{b,a}(s(b)) - s(a)\). When the stalks are
finite-dimensional Hilbert spaces, we equip \(C^0\) and \(C^1\) with
\(\ell^2\) inner products and define the \textbf{sheaf Laplacian}:

\[\Delta_\mathcal{F} = d^* d : C^0(P, F) \to C^0(P, F)\]

where \(d^*\) is the \(\ell^2\)-adjoint of \(d\).

\begin{quote}
\textbf{Theorem 4.1 (Sheaf Laplacian Gradient Structure).}
\textbf{(Layer B)} The sheaf Laplacian \(\Delta_\mathcal{F}\) is a
positive semidefinite operator on \(C^0(P, F)\), and:

\begin{enumerate}
\def\labelenumi{(\alph{enumi})}
\item
  \(\ker(\Delta_\mathcal{F}) = \{s \in C^0(P, F) : \delta_e(s) = 0 \text{ for all } e\}\)
  --- the space of globally consistent sections.
\item
  The gradient flow \(\dot{s} = -\Delta_\mathcal{F} s\) monotonically
  decreases the squared \(\ell^2\) defect
  \(\|ds\|_2^2 = \sum_e \|\delta_e(s)\|^2\) and converges to the nearest
  consistent section.
\item
  The fixed points of the flow are exactly the globally consistent
  sections.
\end{enumerate}
\end{quote}

\emph{Proof.} (a) \(\Delta_\mathcal{F} s = 0\) iff
\(\langle d^*ds, s \rangle = \|ds\|_2^2 = 0\) iff \(ds = 0\) iff
\(\delta_e(s) = 0\) for all \(e\).

\begin{enumerate}
\def\labelenumi{(\alph{enumi})}
\setcounter{enumi}{1}
\item
  Along the flow:
  \(\frac{d}{dt}\|ds\|_2^2 = 2\langle d\dot{s}, ds \rangle = -2\langle d\Delta_\mathcal{F} s, ds \rangle = -2\|d^*ds\|^2 \leq 0\)
  (using the identity \(d \circ d^* d = d d^* d\) and symmetry of
  \(\Delta_\mathcal{F}\)). The flow is a contraction on a
  finite-dimensional space, hence converges.
\item
  Immediate from (a). \(\square\)
\end{enumerate}

\subsubsection{\texorpdfstring{4.2 Relation to \(\ell^1\)
Minimization}{4.2 Relation to \textbackslash ell\^{}1 Minimization}}\label{relation-to-ell1-minimization}

The sheaf Laplacian governs the \(\ell^2\) gradient flow. The \(\ell^1\)
defect minimization studied in Paper 003 is a distinct problem:
minimizing \(\|ds\|_1\) over gauge-equivalent representatives. These are
related but not identical:

\begin{itemize}
\tightlist
\item
  The \(\ell^2\) flow finds the \emph{least-squares} consistent section.
\item
  The \(\ell^1\) quotient norm finds the \emph{sparsest}
  gauge-equivalent representative.
\end{itemize}

As shown in Paper 003, Section 6, the \(\ell^2\) projection generically
inflates \(\ell^1\) cost. The sheaf Laplacian is therefore the natural
operator for characterizing \emph{where} defects live (its spectrum
encodes the obstruction structure), while the \(\ell^1\) quotient norm
measures \emph{how much} irreducible defect remains.

\subsubsection{4.3 Cohomological Selection (Layer C ---
Conjectural)}\label{cohomological-selection-layer-c-conjectural}

The following is stated as a research direction, not a theorem.

\emph{Conjecture 4.2.} In systems where the defect field evolves under a
combination of \(\ell^1\) gradient descent and local stochastic
perturbation, the long-time behavior selects configurations minimizing
the rank of \(H^1\) obstruction. States with lower-dimensional
cohomology classes are dynamically favored.

This conjecture, if established, would connect the structural results of
Papers 000--003 to a dynamical selection principle: physical systems do
not merely \emph{have} obstructions, but \emph{evolve to minimize} them.
The formal development of this idea requires the autopoietic cohomology
framework developed in subsequent work.

\begin{center}\rule{0.5\linewidth}{0.5pt}\end{center}

\subsection{5. The Bridge to Quantum Mechanics (Layer C --- Research
Program)}\label{the-bridge-to-quantum-mechanics-layer-c-research-program}

\subsubsection{5.1 The Derivation Chain}\label{the-derivation-chain}

The results of Papers 000--004 establish: - The \(\ell^1\) coboundary
norm is the unique additive defect diagnostic (Papers 001, 002). -
Projection dynamics generate intrinsic defects measured by this norm
(Paper 000). - The defect has a canonical gauge-invariant cohomological
structure (Paper 003). - The norm is stable under coarse-graining and
extends to operator-valued settings (this paper).

Paper 000 (Appendix A) outlines a derivation chain from these structural
results to the emergence of quantum mechanical features:

\[\text{Observer consistency} \xrightarrow{\S 1.2} \text{Axioms 1--4} \xrightarrow{\text{Thm 4.2}} \ell^1 \xrightarrow{\text{Subsequent}} \text{DFT} \xrightarrow{\text{Parseval}} P = |\psi|^2\]

Each arrow represents a mathematical implication, either proved (solid)
or conjectured (dashed). The present paper strengthens the first three
arrows. The latter steps are developed in subsequent work and are not
established in this paper.

\subsubsection{5.2 Connection to the Fermion Sign
Problem}\label{connection-to-the-fermion-sign-problem}

The Schatten-1 extension (Theorem 3.1) has a direct application to the
fermion sign problem. When the presheaf stalks carry fermionic
statistics, the restriction maps \(r_{b,a}\) involve antisymmetrization,
and the coboundary defect \(\delta_e\) measures the failure of local
fermionic states to glue consistently. The trace norm of this defect is
identified with sign frustration {[}9{]}.

This provides a structural interpretation: the fermion sign problem is
not an artifact of computational representation, but a manifestation of
\(\ell^1\) gauge frustration on operator-valued presheaves.

\subsubsection{5.3 Connection to Coupling
Constants}\label{connection-to-coupling-constants}

Subsequent work {[}10{]} derives the fine structure constant from a
minimal \(\ell^1\) obstruction seed. The key bridge from the present
paper is the observation that:

\begin{enumerate}
\def\labelenumi{\arabic{enumi}.}
\tightlist
\item
  The \(\ell^1\) obstruction on a 2-node graph has coboundary norm 2
  (the minimal non-trivial case).
\item
  Subdivision healing (a specific instance of the coarse-graining
  inverse) produces a 3-node lattice.
\item
  By Theorem 2.1, the defect is monotonically preserved through this
  process.
\item
  The Schatten-1 extension (Theorem 3.1) ensures the framework applies
  to quantum states on this lattice.
\end{enumerate}

The subsequent derivation of wave mechanics, the Born rule, and coupling
constants from this lattice is developed in further work. The present
paper provides the structural justification for why the \(\ell^1\)
framework extends to that setting.

\begin{center}\rule{0.5\linewidth}{0.5pt}\end{center}

\subsection{6. Discussion}\label{discussion}

\subsubsection{6.1 What This Paper Proves}\label{what-this-paper-proves}

Three new results extend the foundational quad:

\begin{enumerate}
\def\labelenumi{\arabic{enumi}.}
\item
  \textbf{Functorial aggregation is monotone} (Theorem 2.1): \(\ell^1\)
  defects cannot be hidden by coarse-graining. This is the
  renormalization stability of the obstruction framework.
\item
  \textbf{The operator-valued extension is canonical} (Theorem 3.1):
  When stalks are operator algebras, the trace norm is the unique
  diagnostic satisfying the same axioms. Quantum marginal
  incompatibility and entanglement monogamy are natural consequences.
\item
  \textbf{The sheaf Laplacian governs defect dynamics} (Theorem 4.1):
  Its kernel is the space of globally consistent sections, and its
  gradient flow monotonically reduces defect.
\end{enumerate}

\subsubsection{6.2 What This Paper Does Not
Prove}\label{what-this-paper-does-not-prove}

The conjectural elements (Layer C) are explicitly marked:

\begin{itemize}
\tightlist
\item
  The dynamical selection principle (Conjecture 4.2) is a research
  direction.
\item
  The derivation chain to quantum mechanics (Section 5.1) contains both
  established and open steps.
\item
  The connections to applications in subsequent work are structural
  observations, not formal proofs.
\end{itemize}

The conditionality of all claims is preserved. The \(\ell^1\) framework
is forced \emph{given} the axioms of Papers 001--002. The extensions
proved here are forced \emph{given} the \(\ell^1\) framework. The
conjectural elements indicate where the program requires further formal
development.

\subsubsection{6.3 Position in the
Program}\label{position-in-the-program}

Papers 000--003 answer: \emph{What must the defect metric be?}
(\(\ell^1\), uniquely.)

This paper answers: \emph{How does it behave?} (Monotone under
aggregation, canonical on operators, governed by the sheaf Laplacian.)

Subsequent work explores: \emph{What follows?} (Autopoietic cohomology,
gauge emergence, sign frustration, coupling constants, grand
unification, cosmological synthesis.)

The program is falsifiable at every joint: if any proved theorem is
incorrect, or if any conjectured step fails, the downstream structure
requires revision at that point. The separation of established results
(Layers A, B) from conjectures (Layer C) is maintained to make this
falsification structure transparent.

\begin{center}\rule{0.5\linewidth}{0.5pt}\end{center}

\subsection{Acknowledgments}\label{acknowledgments}

No external funding. No conflicts of interest.

\begin{center}\rule{0.5\linewidth}{0.5pt}\end{center}

\subsection{References}\label{references}

{[}1{]} S. Mac Lane, \emph{Categories for the Working Mathematician},
Springer, 1971.

{[}2{]} B. Eckmann, ``Harmonische Funktionen und Randwertaufgaben in
einem Komplex,'' \emph{Commentarii Mathematici Helvetici} \textbf{17}
(1944/45), 240--255.

{[}3{]} J. A. Hansen and R. Ghrist, ``Toward a spectral theory of
cellular sheaves,'' \emph{Journal of Applied and Computational Topology}
\textbf{3} (2019), 315--358.

{[}4{]} J. H. Carroll, ``Projection Obstruction Theory: Retraction
Nonlinearity, \(\ell^1\) Rigidity, and Density Scaling,'' Zenodo, 2026.
https://doi.org/10.5281/zenodo.19151803

{[}5{]} J. H. Carroll, ``The Free \(\ell^1\) Seminorm on Banach Presheaf
Coboundaries,'' Zenodo, 2026. https://doi.org/10.5281/zenodo.19152115

{[}6{]} J. H. Carroll, ``Coordinate-Wise Additivity and the \(\ell^1\)
Norm on Finite Graph Cochains,'' Zenodo, 2026.
https://doi.org/10.5281/zenodo.19152599

{[}7{]} J. H. Carroll, ``Hodge Structure and Gauge Equivalence of
\(\ell^1\) Defect Fields,'' Zenodo, 2026.
https://doi.org/10.5281/zenodo.19152799

{[}8{]} R. T. Rockafellar, \emph{Convex Analysis}, Princeton University
Press, 1970.

\end{document}
